\section{Definícia problému a cieľ práce} \label{s: problem}

V reálnej prevádzke musia distribuované systémy spĺňať prísne požiadavky na úroveň poskytovaných služieb (SLA). Hlavným problémom pri návrhu komunikácie je kompromis (trade-off) medzi rýchlosťou prenosu dát a garanciou ich doručenia.

Tradičné synchrónne prístupy založené na protokole HTTP/REST a formáte JSON trpia vysokou réžiou dát a blokovaním vlákien. Moderné alternatívy ako gRPC využívajú binárnu serializáciu (Protocol Buffers) a multiplexovanie cez HTTP/2, čo výrazne znižuje latenciu~\cite{Indrasiri2020}. Na druhej strane, asynchrónna komunikácia pomocou brokerov správ (Apache Kafka, NATS JetStream) oddeľuje odosielateľa od príjemcu, čím zvyšuje odolnosť voči dočasným výpadkom.

Cieľom tohto projektu je kvantifikovať správanie týchto technológií pri umelo vyvolanom sieťovom chaose a posúdiť ich vhodnosť pre kritické distribuované aplikácie.